% Proof of the generalised spider lemma: the three degenerate cases, then the
% 2x2 core that drives the induction.
\begin{center}\begin{tabular}{r c c c}
  \diagrow{$a=0$}{
      \pic (g1) at (0.500,0.500) {codiscard};
      \pic (g2) at (1.750,0.500) {copy};
      \pic (g3) at (3.000,1.000) {ge};
      \pic (g4) at (3.000,0.000) {ge};
      \draw (g1-out-0) -- (g2-in-0);
      \draw (g2-out-0) to[out=0,in=180] (g3-in-0);
      \draw (g2-out-1) to[out=0,in=180] (g4-in-0);
  }{
      \pic (g1) at (0.500,1.000) {codiscard};
      \pic (g2) at (0.500,0.000) {codiscard};
      \draw (g1-out-0) -- (1.250,1.000);
      \draw (g2-out-0) -- (1.250,0.000);
  }
  \diagrow{$b=0$}{
      \pic (g1) at (0.500,1.000) {ge};
      \pic (g2) at (0.500,0.000) {ge};
      \pic (g3) at (1.750,0.500) {cocopy};
      \pic (g4) at (3.000,0.500) {discard};
      \draw (g1-out-0) to[out=0,in=180] (g3-in-0);
      \draw (g2-out-0) to[out=0,in=180] (g3-in-1);
      \draw (g3-out-0) -- (g4-in-0);
  }{
      \pic (g1) at (0.500,1.000) {discard};
      \pic (g2) at (0.500,0.000) {discard};
      \draw (0.000,1.000) -- (g1-in-0);
      \draw (0.000,0.000) -- (g2-in-0);
  }
  \diagrow{$a=1$}{
      \pic (g1) at (0.500,0.500) {ge};
      \pic (g2) at (1.750,0.500) {copy};
      \pic (g3) at (3.000,1.000) {ge};
      \pic (g4) at (3.000,0.000) {ge};
      \draw (g1-out-0) to[out=0,in=180] (g2-in-0);
      \draw (g2-out-0) to[out=0,in=180] (g3-in-0);
      \draw (g2-out-1) to[out=0,in=180] (g4-in-0);
  }{
      \pic (g1) at (0.500,0.500) {copy};
      \pic (g2) at (1.750,1.000) {ge};
      \pic (g3) at (1.750,0.000) {ge};
      \draw (g1-out-0) to[out=0,in=180] (g2-in-0);
      \draw (g1-out-1) to[out=0,in=180] (g3-in-0);
  }
  \diagrow{$b=1$}{
      \pic (g1) at (0.500,1.000) {ge};
      \pic (g2) at (0.500,0.000) {ge};
      \pic (g3) at (1.750,0.500) {cocopy};
      \pic (g4) at (3.000,0.500) {ge};
      \draw (g1-out-0) to[out=0,in=180] (g3-in-0);
      \draw (g2-out-0) to[out=0,in=180] (g3-in-1);
      \draw (g3-out-0) to[out=0,in=180] (g4-in-0);
  }{
      \pic (g1) at (0.500,1.000) {ge};
      \pic (g2) at (0.500,0.000) {ge};
      \pic (g3) at (1.750,0.500) {cocopy};
      \draw (g1-out-0) to[out=0,in=180] (g3-in-0);
      \draw (g2-out-0) to[out=0,in=180] (g3-in-1);
  }
\end{tabular}\end{center}

\vspace{2mm}

\begin{center}
    \begin{tikzpicture}[axpic]
        \pic (g1) at (0.500,1.000) {ge};
        \pic (g2) at (0.500,0.000) {ge};
        \pic (c1) at (1.750,0.500) {cocopy};
        \pic (c2) at (3.000,0.500) {copy};
        \pic (h1) at (4.250,1.000) {ge};
        \pic (h2) at (4.250,0.000) {ge};
        \draw (g1-out-0) to[out=0,in=180] (c1-in-0);
        \draw (g2-out-0) to[out=0,in=180] (c1-in-1);
        \draw (c1-out-0) -- (c2-in-0);
        \draw (c2-out-0) to[out=0,in=180] (h1-in-0);
        \draw (c2-out-1) to[out=0,in=180] (h2-in-0);
    \end{tikzpicture}
    \;\begin{tikzpicture}[axpic]\node{$\overset{(\text{G4}),(e),(f)}{=}$};\end{tikzpicture}\;
    \begin{tikzpicture}[axpic]
        \pic (c1) at (0.500,1.500) {copy};
        \pic (c2) at (0.500,-0.500) {copy};
        \pic (g11) at (2.000,2.000) {ge};
        \pic (g12) at (2.000,1.000) {ge};
        \pic (g21) at (2.000,0.000) {ge};
        \pic (g22) at (2.000,-1.000) {ge};
        \pic (d1) at (3.750,1.500) {cocopy};
        \pic (d2) at (3.750,-0.500) {cocopy};
        \draw (c1-out-0) to[out=0,in=180] (g11-in-0);
        \draw (c1-out-1) to[out=0,in=180] (g12-in-0);
        \draw (c2-out-0) to[out=0,in=180] (g21-in-0);
        \draw (c2-out-1) to[out=0,in=180] (g22-in-0);
        \draw (g11-out-0) to[out=0,in=180] (d1-in-0);
        \draw (g21-out-0) to[out=0,in=180] (d1-in-1);
        \draw (g12-out-0) to[out=0,in=180] (d2-in-0);
        \draw (g22-out-0) to[out=0,in=180] (d2-in-1);
    \end{tikzpicture}
\end{center}
